\documentclass[12pt,aspectratio=169,ignorenonframetext]{ltxbeamer}

\usepackage[spanish,es-nodecimaldot,es-tabla]{babel}

\title{Elementos Flotantes}
\subtitle{Curso de {\lmr\LaTeX}}
\author{Joar Buitrago\texorpdfstring{\thanks{jebuitragoc@unal.edu.co}}{ <jebuitragoc@unal.edu.co>}}
\date{}

\usepackage{verbatim}

\usetikzlibrary{spy}

\addbibresource{bibliography.bib}



\begin{document}
\begin{frame}
\maketitle
\end{frame}





\begin{frame}
\frametitle{El Paquete \texttt{graphicx}}

El estándar de facto para incluir contenido gráfico en documentos {\lmr\LaTeX} es el paquete \texttt{graphicx}.

\bigskip\pause

\begin{columns}
    \begin{column}{0.2\textwidth}
        {\Large Raster}

        \begin{itemize}
            \item JPEG
            \item PNG
            \item BMP
        \end{itemize}
    \end{column}
    \begin{column}{0.2\textwidth}
        {\Large Vector}

        \begin{itemize}
            \item EPS
            \item PS
            \item PDF
        \end{itemize}
    \end{column}
\end{columns}
\end{frame}





\section{El paquete \texttt{graphicx}}

% Muchos de los ejemplos que verán en esta sección no corresponden con lo que se muestra como su código, esto es porque las imágenes de muestra son muy grandes, y LaTeX las colocará acorde a su densidad o PPP (Pixeles Por Pulgada). La mayoría de imágenes tienen una densidad de 72 PPP (de aquí la definición de `big point` `bp`).

% Por ejemplo, la imagen `assets/beach.jpg` tiene 1200 * 800 px, y una densidad de 72 PPP, que haciendo un cálculo sencillo podemos notar rápidamente que la imagen tiene un tamaño físico de 42 1/3 * 28 2/9 cm, que es muy superior al tamaño de una sola diapositiva de esta presentación (que a una relación de aspecto 16:9 tiene 16 * 9 cm).

\begin{frame}
\frametitle{El Paquete \texttt{graphicx}}
\begin{columns}
    \begin{column}{0.3\textwidth}
        \centering
        {\Large Raster}\bigskip

        \begin{tikzpicture}[
            spy using outlines={
                rectangle,
                red,
                magnification=25,
                size=2cm,
                connect spies,
            },
        ]
            \node [inner sep=0] (element) {\includegraphics{assets/A.png}};

            \draw (element.north) ++(0.05,-0.15) coordinate (spy);

            \spy [] on (spy) in node [right of=element,xshift=1cm];
        \end{tikzpicture}
    \end{column}
    \begin{column}{0.3\textwidth}
        \centering
        {\Large Vector}\bigskip

        \begin{tikzpicture}[
            spy using outlines={
                rectangle,
                red,
                magnification=25,
                size=2cm,
                connect spies,
            },
        ]
            \node [inner sep=0] (element) {\includegraphics{assets/A.pdf}};

            \draw (element.north) ++(0.05,-0.15) coordinate (spy);

            \spy [] on (spy) in node [right of=element,xshift=1cm];
        \end{tikzpicture}
    \end{column}
\end{columns}
\end{frame}





\begin{frame}[fragile]
\begin{minted}[escapeinside=||]{latex}
\includegraphics[<|\textit{opciones clave=valor}|>]{<|\textit{archivo}|>}
\end{minted}
\end{frame}

\begin{frame}[fragile]
\begin{columns}
    \begin{column}{0.5\textwidth}
        \begin{minted}{latex}
\includegraphics{assets/beach.jpg}
        \end{minted}
    \end{column}
    {\color{gray}\vline}
    \begin{column}{0.5\textwidth}
        \lmr
        \includegraphics[width=0.95\textwidth]{assets/beach.jpg}
    \end{column}
\end{columns}
\end{frame}





\begin{frame}[fragile]
\begin{columns}
    \begin{column}{0.5\textwidth}
        \begin{minted}{latex}
\includegraphics{
    C:/Users/User/Images/beach.jpg
}
        \end{minted}
    \end{column}
    {\color{gray}\vline}
    \begin{column}{0.5\textwidth}
        \lmr
        \includegraphics[width=0.95\textwidth]{assets/beach.jpg}
    \end{column}
\end{columns}
\end{frame}





\begin{frame}
\begin{alertblock}{IMPORTANTE}
    Todas las gráficas incluidas por {\lmr\LaTeX} son \alert{cajas}.
\end{alertblock}
\end{frame}





\begin{frame}[fragile]
\frametitle{Opciones de inclusión de gráficas}
\framesubtitle{\texttt{scale}}

\texttt{scale=\textit{escala}} Escala la imagen proporcionalmente.

\bigskip\pause

\begin{columns}
    \begin{column}{0.5\textwidth}
        \begin{minted}{latex}
\includegraphics[
    scale=1,
]{assets/beach.jpg}

\includegraphics[
    scale=0.5,
]{assets/beach.jpg}
        \end{minted}
    \end{column}
    {\color{gray}\vline}
    \begin{column}{0.5\textwidth}
        \lmr
        \includegraphics[width=0.8\textwidth]{assets/beach.jpg}

        \includegraphics[width=0.4\textwidth]{assets/beach.jpg}
    \end{column}
\end{columns}
\end{frame}




\begin{frame}[fragile]
\frametitle{Opciones de inclusión de gráficas}
\framesubtitle{\texttt{width} y \texttt{height}}

\texttt{width=\textit{ancho}} y \texttt{height=\textit{alto}} definen el tamaño de la imagen. Reciben una dimensión de {\lmr\TeX}.

\bigskip\pause

\begin{columns}
    \begin{column}{0.5\textwidth}
        \begin{minted}{latex}
\includegraphics[
    width=3cm,
    height=4cm,
]{assets/beach.jpg}
        \end{minted}
    \end{column}
    {\color{gray}\vline}
    \begin{column}{0.5\textwidth}
        \lmr
        \includegraphics[width=3cm,height=4cm]{assets/beach.jpg}
    \end{column}
\end{columns}
\end{frame}





\begin{frame}[fragile]
\frametitle{Opciones de inclusión de gráficas}
\framesubtitle{\texttt{width} y \texttt{height}}

Podemos definir solamente una de las dos y la otra será calculada de manera automática.

\bigskip

\begin{columns}
    \begin{column}{0.5\textwidth}
        \begin{minted}{latex}
\includegraphics[
    width=3cm,
]{assets/beach.jpg}
        \end{minted}
    \end{column}
    {\color{gray}\vline}
    \begin{column}{0.5\textwidth}
        \lmr
        \includegraphics[width=3cm]{assets/beach.jpg}
    \end{column}
\end{columns}
\end{frame}





\begin{frame}[fragile]
\frametitle{Opciones de inclusión de gráficas}
\framesubtitle{\texttt{width} y \texttt{height}}

{\lmr\LaTeX} posee comandos especiales que nos dan longitudes dependiendo de nuestro contexto.

\bigskip\pause

\begin{center}
\scriptsize
\begin{tabular}{>{\ttfamily}ll}
    \toprule
    {\sffamily Comando} & Definición \\
    \midrule
    \mintinline{latex}{\linewidth} & Ancho de la línea en el entorno actual \\
    \mintinline{latex}{\textwidth} & Ancho del área de texto en la página \\
    \mintinline{latex}{\textheight} & Altura del área de texto en la página \\
    \mintinline{latex}{\paperwidth} & Ancho de la página \\
    \mintinline{latex}{\paperheight} & Altura de la página \\
    \mintinline{latex}{\columnwidth} & Ancho de la columna \\
    \mintinline{latex}{\columnsep} & Separación entre columnas \\
    \bottomrule
\end{tabular}
\end{center}
\end{frame}



\begin{frame}[fragile]
\frametitle{Opciones de inclusión de gráficas}
\framesubtitle{\texttt{width} y \texttt{height}}

Nosotros podemos definir nuestras propias longitudes o modificar otras ya definidas por {\lmr\LaTeX}.

\bigskip

\begin{minted}[escapeinside=||]{latex}
\newlength{|\textit{\textbackslash nombre}|}
\setlength{|\textit{\textbackslash nombre}|}{|\textit{dimensión}|}
\end{minted}
\end{frame}






\begin{frame}[fragile]
\frametitle{Opciones de inclusión de gráficas}
\framesubtitle{\texttt{draft}}

\texttt{draft} cambia la inclusión de la imagen a modo borrador. Muy útil cuando la imagen es muy pesada, y se desea excluir temporalmente de la compilación.

\bigskip\pause

\begin{columns}
    \begin{column}{0.5\textwidth}
        \begin{minted}{latex}
\includegraphics[
    draft,
]{assets/beach.jpg}
        \end{minted}
    \end{column}
    {\color{gray}\vline}
    \begin{column}{0.5\textwidth}
        \lmr
        \includegraphics[width=0.95\linewidth,draft]{assets/beach.jpg}
    \end{column}
\end{columns}
\end{frame}





\begin{frame}[fragile]
\frametitle{Opciones de inclusión de gráficas}
\framesubtitle{\texttt{angle} y \texttt{origin}}

\texttt{angle=\textit{ángulo}} y \texttt{origin=\textit{posición}} definen la rotación de la gráfica incluida.

\bigskip\pause

\begin{columns}
    \begin{column}{0.5\textwidth}
        \begin{minted}{latex}
\includegraphics[
    angle=30,
    origin=c,
]{assets/beach.jpg}
        \end{minted}
    \end{column}
    {\color{gray}\vline}
    \begin{column}{0.5\textwidth}
        \lmr
        \includegraphics[width=0.5\linewidth,angle=30,origin=c]{assets/beach.jpg}
    \end{column}
\end{columns}
\end{frame}





\begin{frame}[fragile]
\frametitle{Opciones de inclusión de gráficas}
\framesubtitle{\texttt{angle} y \texttt{origin}}
\begin{columns}
    \begin{column}{0.5\textwidth}
        \begin{minted}{latex}
texto
\includegraphics[
    angle=180,
    origin=c,
]{assets/beach.jpg}
texto
        \end{minted}
    \end{column}
    {\color{gray}\vline}
    \begin{column}{0.5\textwidth}
        \lmr
        texto
        \includegraphics[width=0.5\linewidth,angle=180,origin=c]{assets/beach.jpg}
        texto
    \end{column}
\end{columns}
\end{frame}





\begin{frame}[fragile]
\frametitle{Opciones de inclusión de gráficas}
\framesubtitle{\texttt{angle} y \texttt{origin}}
\begin{columns}
    \begin{column}{0.5\textwidth}
        \begin{minted}{latex}
texto
\includegraphics[
    angle=180,
    origin=t,
]{assets/beach.jpg}
texto
        \end{minted}
    \end{column}
    {\color{gray}\vline}
    \begin{column}{0.5\textwidth}
        \lmr
        texto
        \includegraphics[width=0.5\linewidth,angle=180,origin=t]{assets/beach.jpg}
        texto
    \end{column}
\end{columns}
\end{frame}



\begin{frame}[fragile]
\frametitle{Opciones de inclusión de gráficas}
\framesubtitle{\texttt{angle} y \texttt{origin}}
\begin{columns}
    \begin{column}{0.5\textwidth}
        \begin{minted}{latex}
texto
\includegraphics[
    angle=180,
    origin=b,
]{assets/beach.jpg}
texto
        \end{minted}
    \end{column}
    {\color{gray}\vline}
    \begin{column}{0.5\textwidth}
        \lmr
        texto
        \includegraphics[width=0.5\linewidth,angle=180,origin=b]{assets/beach.jpg}
        texto
    \end{column}
\end{columns}
\end{frame}





\section{Entornos Flotantes}

\begin{frame}[fragile]
\frametitle{¿Qué son los Entornos Flotantes?}

Son bloques de contenido que {\lmr\LaTeX} decide dónde colocar mejor para optimizar el diseño de la página.

\bigskip\pause

\begin{itemize}[<+->]
    \item Asignan un número secuencial (e.g., Figura 1, Tabla 2).
    \item Permiten añadir una leyenda con \mintinline{latex}{\caption}.
    \item Permiten etiquetar el entorno usando \mintinline{latex}{\label} para hacer referencias cruzadas.
\end{itemize}
\end{frame}





\begin{frame}
\frametitle{Especificadores de Posicionamiento}

Los entornos flotantes toman un argumento opcional que sugiere la ubicación

\begin{center}
    \scriptsize
    \begin{tabular}{>{\ttfamily}lll}
        \toprule
        {\sffamily Comando} & Significado & Acción \\
        \midrule
        h & here & Ubicar en donde se utiliza el entorno. \\
        t & top & Colocarlo en la parte \textbf{superior} de una página. \\
        b & bottom & Colocarlo en la parte \textbf{inferior} de una página. \\
        p & page & Colocarlo en una \textbf{página separada} de flotantes. \\
        ! & override & Ignora la mayoría de los parámetros internos de {\lmr\LaTeX}. \\
        \bottomrule
    \end{tabular}
\end{center}

\pause

Por defecto {\lmr\LaTeX} utiliza el especificador \texttt{tbp}.
\end{frame}




\begin{frame}
\begin{alertblock}{IMPORTANTE}
    El especificador de ubicación es solamente una \textbf{sugerencia} que hacemos a {\lmr\LaTeX}.

    \bigskip

    {\lmr\LaTeX} internamente utilizará diversos parámetros para acomodarse a esa sugerencia, pero si no encuentra como ubicar el contenido del entorno dentro de esos parámetros, la ignorará.
\end{alertblock}
\end{frame}




\begin{frame}[fragile]
\frametitle{Referencias Cruzadas}

Para añadir numeración y leyenda a los elementos flotantes utilizaremos el comando \mintinline{latex}{\caption}.

\pause

\begin{minted}[escapeinside=||]{latex}
\caption[|\textit{short title}|]{|\textit{title}|}
\end{minted}
\end{frame}






\begin{frame}[fragile]
\frametitle{Referencias Cruzadas}

Para etiquetar la numeración que asigno {\lmr\LaTeX} al entorno flotante, utilizaremos \mintinline{latex}{\label} dentro o después del comando \mintinline{latex}{\caption}

\bigskip\pause

\begin{columns}
    \begin{column}{0.5\textwidth}
        \begin{minted}[escapeinside=||]{latex}
\caption{Cuerpo A \label{flt:b1}}

\caption{Cuerpo B}
\label{flt:b2}

\ref{flt:b1} \ref{flt:b2}
        \end{minted}
    \end{column}
    {\color{gray}\vline}
    \begin{column}{0.5\textwidth}
        \lmr
        \begingroup
        \centering
        Flotante 1: Cuerpo A\par
        \endgroup

        \begingroup
        \centering
        Flotante 2: Cuerpo B\par
        \endgroup

        1 2
    \end{column}
\end{columns}
\end{frame}





\begin{frame}
\frametitle{Entornos Flotantes}

Los dos entornos flotantes que incluye {\lmr\LaTeX} por defecto son:

\begin{description}
    \item[\texttt{figure}] Para figuras, gráficas o demás elementos visuales.
    \item[\texttt{table}] Para tablas.
\end{description}

\pause

Cada uno de estos entornos incluye una versión estrellada que coloca sus contenidos a lo ancho de toda la página en caso de que estemos creando un documento a dos columnas.
\end{frame}





\begin{frame}[fragile]
\frametitle{Entornos Flotantes}

\begin{columns}
    \begin{column}{0.5\textwidth}
        \begin{minted}[escapeinside=||]{latex}
\begin{figure}[htb]
    \includegraphics[
        width=0.5\linewidth,
    ]{assets/beach.jpg}
    \caption{Una hermosa playa}
    \label{fig:playa}
\end{figure}
        \end{minted}
    \end{column}
    \pause
    {\color{gray}\vline}
    \begin{column}{0.5\textwidth}
        \lmr
        \includegraphics[
            width=0.5\linewidth,
        ]{assets/beach.jpg}

        \begingroup
        \centering
        Figura 1: Una hermosa playa\par
        \endgroup
    \end{column}
\end{columns}
\end{frame}






\begin{frame}[fragile]
\frametitle{Alineación}

Para realizar la alineación del texto, {\lmr\LaTeX} ofrece diversas opciones para este fin.

\pause\bigskip

\begin{columns}
    \begin{column}{0.3\textwidth}
        {\Large Izquierda}\bigskip

        \begin{itemize}
            \item \mintinline{latex}{\raggedright}
            \item \texttt{flushleft}
        \end{itemize}
    \end{column}
    \begin{column}{0.3\textwidth}
        {\Large Centro}\bigskip

        \begin{itemize}
            \item \mintinline{latex}{\centering}
            \item \texttt{center}
        \end{itemize}
    \end{column}
    \begin{column}{0.3\textwidth}
        {\Large Derecha}\bigskip

        \begin{itemize}
            \item \mintinline{latex}{\raggedleft}
            \item \texttt{flushright}
        \end{itemize}
    \end{column}
\end{columns}

\pause\bigskip

\begin{block}{Entorno Caprichoso}
    El entorno \texttt{center} añade espacio extra encima y debajo de su contenido.
\end{block}
\end{frame}






\begin{frame}[fragile]
\frametitle{Entornos Flotantes}

\begin{columns}
    \begin{column}{0.5\textwidth}
        \begin{minted}[escapeinside=||]{latex}
\begin{figure}[htb]
    \includegraphics[
        width=0.5\linewidth,
    ]{assets/beach.jpg}
    \caption{Una hermosa playa}
    \label{fig:playa}
\end{figure}
        \end{minted}
    \end{column}
    {\color{gray}\vline}
    \begin{column}{0.5\textwidth}
        \lmr
        \includegraphics[
            width=0.5\linewidth,
        ]{assets/beach.jpg}

        \medskip
        \begingroup
        \centering
        Figura 1: Una hermosa playa\par
        \endgroup
    \end{column}
\end{columns}
\end{frame}






\begin{frame}[fragile]
\frametitle{Entornos Flotantes}

\begin{columns}
    \begin{column}{0.5\textwidth}
        \begin{minted}[escapeinside=||]{latex}
\begin{figure}[htb]
    \centering
    \includegraphics[
        width=0.5\linewidth,
    ]{assets/beach.jpg}
    \caption{Una hermosa playa}
    \label{fig:playa}
\end{figure}
        \end{minted}
    \end{column}
    {\color{gray}\vline}
    \begin{column}{0.5\textwidth}
        \lmr
        \centering
        \includegraphics[
            width=0.5\linewidth,
        ]{assets/beach.jpg}

        \medskip
        \begingroup
        \centering
        Figura 1: Una hermosa playa\par
        \endgroup
    \end{column}
\end{columns}
\end{frame}





\begin{frame}[fragile]
\frametitle{Entornos Flotantes}

\begin{columns}
    \begin{column}{0.5\textwidth}
        \begin{minted}[escapeinside=||]{latex}
\begin{table}[htb]
    \begin{tabular}{ccc}
        \toprule
        Given Name & Surname \\
        \midrule
        Biggus & Dickus \\
        Monty & Python \\
        \bottomrule
    \end{tabular}
    \caption{Listado de usuarios}
    \label{tab:usuarios}
\end{figure}
        \end{minted}
    \end{column}
    \pause
    {\color{gray}\vline}
    \begin{column}{0.5\textwidth}
        \lmr
        \begin{tabular}{cc}
            \toprule
            Given Name & Surname \\
            \midrule
            Biggus & Dickus \\
            Monty & Python \\
            \bottomrule
        \end{tabular}

        \medskip
        \begingroup
        \centering
        Cuadro 1: Listado de usuarios\par
        \endgroup
    \end{column}
\end{columns}
\end{frame}






\begin{frame}[fragile]
\frametitle{Entornos Flotantes}

\begin{columns}
    \begin{column}{0.5\textwidth}
        \begin{minted}[escapeinside=||]{latex}
\begin{table}[htb]
    \centering
    \begin{tabular}{ccc}
        \toprule
        Given Name & Surname \\
        \midrule
        Biggus & Dickus \\
        Monty & Python \\
        \bottomrule
    \end{tabular}
    \caption{Listado de usuarios}
    \label{tab:usuarios}
\end{figure}
        \end{minted}
    \end{column}
    {\color{gray}\vline}
    \begin{column}{0.5\textwidth}
        \lmr
        \centering
        \begin{tabular}{cc}
            \toprule
            Given Name & Surname \\
            \midrule
            Biggus & Dickus \\
            Monty & Python \\
            \bottomrule
        \end{tabular}

        \medskip
        \begingroup
        \centering
        Cuadro 1: Listado de usuarios\par
        \endgroup
    \end{column}
\end{columns}
\end{frame}







\section{El paquete \texttt{float}}

\begin{frame}
\frametitle{El Paquete \texttt{float}}

Al momento de trabajar con elementos flotantes, puede llegar a ser frustrante el no tener control sobre la posición de los entornos flotantes.

\pause

El paquete \texttt{float} nos ofrece una herramienta para lidiar con esto.
\end{frame}





\begin{frame}
\frametitle{El especificador \texttt{H}}

El especificador de posición \texttt{H} que incluye el paquete \texttt{float} fuerza a {\lmr\LaTeX} a ubicar el entorno flotante en la posición en la que es utilizado en el documento.
\end{frame}






\begin{frame}[fragile]
\frametitle{El especificador \texttt{H}}

\begin{columns}
    \begin{column}{0.5\textwidth}
        \begin{minted}[escapeinside=||]{latex}
\begin{table}[H]
    \centering
    \begin{tabular}{ccc}
        \toprule
        Given Name & Surname \\
        \midrule
        Biggus & Dickus \\
        Monty & Python \\
        \bottomrule
    \end{tabular}
    \caption{Listado de usuarios}
    \label{tab:usuarios}
\end{figure}
        \end{minted}
    \end{column}
    {\color{gray}\vline}
    \begin{column}{0.5\textwidth}
        \lmr
        \centering
        \begin{tabular}{cc}
            \toprule
            Given Name & Surname \\
            \midrule
            Biggus & Dickus \\
            Monty & Python \\
            \bottomrule
        \end{tabular}

        \medskip
        \begingroup
        \centering
        Cuadro 1: Listado de usuarios\par
        \endgroup
    \end{column}
\end{columns}
\end{frame}





\begin{frame}[fragile]
\frametitle{Entornos Flotantes}

\begin{columns}
    \begin{column}{0.5\textwidth}
        \begin{minted}[escapeinside=||]{latex}
\begin{figure}[htb]
    \centering
    \includegraphics[
        width=0.5\linewidth,
    ]{assets/beach.jpg}
    \caption{Una hermosa playa}
    \label{fig:playa}
\end{figure}
        \end{minted}
    \end{column}
    {\color{gray}\vline}
    \begin{column}{0.5\textwidth}
        \lmr
        \centering
        \includegraphics[
            width=0.5\linewidth,
        ]{assets/beach.jpg}

        \medskip
        \begingroup
        \centering
        Figura 1: Una hermosa playa\par
        \endgroup
    \end{column}
\end{columns}
\end{frame}





\begin{frame}[fragile]
\frametitle{Entornos Flotantes}

\begin{columns}
    \begin{column}{0.5\textwidth}
        \begin{minted}[escapeinside=||]{latex}
\begin{figure}[htb]
    \centering
    \caption{Una hermosa playa}
    \includegraphics[
        width=0.5\linewidth,
    ]{assets/beach.jpg}
    \label{fig:playa}
\end{figure}
        \end{minted}
    \end{column}
    {\color{gray}\vline}
    \begin{column}{0.5\textwidth}
        \lmr
        \centering
        \begingroup
        \centering
        Figura 1: Una hermosa playa\par
        \endgroup

        \medskip
        \includegraphics[
            width=0.5\linewidth,
        ]{assets/beach.jpg}
    \end{column}
\end{columns}
\end{frame}







\begin{frame}[allowframebreaks]
\frametitle{Referencias}
\nocite{*}
\printbibliography
\end{frame}
\end{document}
